\noindent\textbf{Elección del procedimiento.} Para régimen turbulento se usa la correlación de Haaland dada en el capítulo: $f_D=[-1{,}8\log_{10}((\varepsilon/(3.7D))^{1.11}+6.9/\mathrm{Re})]^{-2}$. Primero se calcula Reynolds y después se introduce el factor de Darcy en la pérdida de presión.

\begin{enumerate}

\item La viscosidad es cinemática, por lo que no se añade densidad en Reynolds.

\[\fitmath{\mathrm{Re}=\frac{V\cdot D}{\nu}}\]

\[\fitmath{\mathrm{Re}=\frac{(1)\cdot (0{,}1)}{(1{,}007\times10^{-6})}\approx 99304{,}87\,}\]

\item Usar logaritmo decimal y rugosidad en las mismas unidades que diámetro.

\[\fitmath{\mathrm{fD}={\left(-1{,}8\cdot \frac{\ln\left({\left(\frac{\mathrm{eps}}{3.7\cdot D}\right)}^{1.11}+\frac{6.9}{\mathrm{Re}}\right)}{\ln\left(10\right)}\right)}^{-2}}\]

\[\fitmath{\mathrm{fD}={\left(-1{,}8\cdot \frac{\ln\left({\left(\frac{(4{,}5\times10^{-5})}{3.7\cdot (0{,}1)}\right)}^{1.11}+\frac{6.9}{(99304{,}86594)}\right)}{\ln\left(10\right)}\right)}^{-2}\approx 0{,}01987397\,}\]

\item Aplicar Darcy--Weisbach con el factor de Darcy calculado.

\[\fitmath{\mathrm{dp}=\frac{\frac{\mathrm{fD}\cdot L}{D}\cdot \rho\cdot {V}^{2}}{2}}\]

\[\fitmath{\mathrm{dp}=\frac{\frac{(0{,}01987397315)\cdot (100)}{(0{,}1)}\cdot (998{,}2)\cdot {(1)}^{2}}{2}\approx 9919{,}1\,\mathrm{Pa}}\]

\end{enumerate}
